\section{Large-scale Training Adaptation on Hygon DCU K100}

% \begin{wraptable}{r}{0.52\textwidth}
%   \centering
%   \vspace{-12pt}
%   \caption{Hardware specifications of the Sugon K100 AI accelerator and the training cluster.}
%   \label{tab:k100_specs}
%   \small
%   \renewcommand{\arraystretch}{1.2}
%   \begin{tabular}{ll}
%       \toprule
%       \textbf{Specification} & \textbf{Value} \\
%       \midrule
%       Accelerator & Hygon DCU K100 AI \\
%       FP16/BF16 Performance & 196 TFLOPS \\
%       HBM3 Capacity & 64 GB \\
%       Memory Bandwidth & 896 GB/s \\
%       Host Interface & PCIe 5.0 x16 \\
%       Software Ecosystem & ROCm \\
%       \midrule
%       Nodes in Cluster & 56 \\
%       Accelerators per Node & 4 \\
%       Total Accelerators & 224 \\
%       Inter-node Interconnect & InfiniBand (RDMA) \\
%       \bottomrule
%   \end{tabular}
%   \renewcommand{\arraystretch}{1.0}
%   \vspace{-4pt}
% \end{wraptable}

% \begin{table}[t]
%   \centering
%   \caption{Hardware specifications of the Sugon K100 AI accelerator and the training cluster.}
%   \label{tab:k100_specs}
%   \small
%   \renewcommand{\arraystretch}{1.2}
%   \begin{tabular}{ll}
%       \toprule
%       \textbf{Specification} & \textbf{Value} \\
%       \midrule
%       Accelerator & Hygon DCU K100 AI \\
%       FP16/BF16 Performance & 196 TFLOPS \\
%       HBM3 Capacity & 64 GB \\
%       Memory Bandwidth & 896 GB/s \\
%       Host Interface & PCIe 5.0 x16 \\
%       Software Ecosystem & ROCm \\
%       \midrule
%       Nodes in Cluster & 56 \\
%       Accelerators per Node & 4 \\
%       Total Accelerators & 224 \\
%       Inter-node Interconnect & InfiniBand (RDMA) \\
%       \bottomrule
%   \end{tabular}
%   \renewcommand{\arraystretch}{1.0}
% \end{table}

% The hardware specifications of the Sugon K100 AI accelerator and the 56-node training cluster are summarized in Tab.~\ref{tab:k100_specs}. To ensure long-term sustainability and reduce dependence on foreign hardware supply chains, we deployed the entire training pipeline of JuZhou 1.0 on domestic Sugon K100 AI accelerators (Hygon DCU K100 AI). Each K100 AI unit delivers 196\,TFLOPS of FP16/BF16 half-precision compute, 64\,GB of HBM3 memory, 896\,GB/s memory bandwidth, and PCIe 5.0 x16 connectivity. The training cluster comprises 56 nodes with 4 accelerators per node, totaling 224 K100 AI devices. Built on the ROCm ecosystem, the K100 AI provides deep compatibility with CUDA-based applications and enables a smooth migration path for existing deep learning workloads.
The hardware specifications of the Sugon K100 AI accelerator and the main 56-node training configuration are summarized in Tab.~\ref{tab:k100_specs}. To reduce dependence on foreign hardware supply chains, we deployed the JuZhou~1.0 model training and distillation pipeline on domestic Sugon K100 AI accelerators (Hygon DCU K100 AI). Each K100 AI unit provides 196\,TFLOPS of FP16/BF16 compute, 64\,GB of HBM3 memory, 896\,GB/s memory bandwidth, and PCIe 5.0 x16 connectivity. The main training configuration uses 56 nodes with 4 accelerators per node, totaling 224 K100 AI devices. Built on a ROCm-compatible software stack, the platform supports PyTorch-based diffusion training after targeted runtime adaptation and operator-level optimization.

% Adapting the training pipeline to the K100 AI platform required substantial software-level engineering, carried out in close collaboration with Sugon's engineering team. Starting from the vendor-provided PyTorch Docker image, we upgraded the environment to support the latest versions of the \texttt{transformers} and \texttt{accelerate} libraries, which are essential for modern diffusion model training. A critical challenge was ensuring correct gradient backpropagation through the Rectified Flow formulation on the K100 AI: several operators exhibited numerical discrepancies or unsupported backward passes under the initial runtime. Through targeted debugging and operator-level patching, we resolved these issues and validated gradient correctness end-to-end. Additionally, we implemented custom operator optimizations for performance-critical modules---including attention variants and separable convolutions---to better exploit the K100 AI's hardware characteristics. Collectively, these adaptations established a stable and performant training environment on domestic accelerators for large-scale generative model training.

% Scaling to 56 nodes introduced significant communication bottlenecks across machine boundaries. To address this, we leveraged the InfiniBand (IB) networking modules integrated into the Sugon servers. IB provides high-bandwidth, low-latency interconnects with RDMA (Remote Direct Memory Access) support, enabling direct accelerator-to-accelerator memory access across nodes without CPU involvement. This capability is essential for All-Reduce operations in distributed data-parallel training and proved instrumental in maintaining near-linear scaling efficiency as the cluster grew. The combination of IB interconnects and the feature pre-computation strategy described in Section~4 enabled the 224-device cluster to operate at high throughput throughout both low-resolution pre-training and progressive resolution scaling.

Adapting the training pipeline to the K100 AI platform required substantial software engineering in collaboration with Sugon's engineering team. Starting from the vendor-provided PyTorch Docker image, we upgraded the runtime to support modern diffusion training dependencies, including \texttt{transformers} and \texttt{accelerate}. A key challenge was ensuring correct gradient backpropagation under the Rectified Flow formulation, as several operators showed numerical discrepancies or unsupported backward passes in the initial environment. Through targeted debugging, operator-level patching, and end-to-end gradient validation, we established a stable training runtime. We further optimized performance-critical modules, including attention variants and separable convolutions, to better match the K100 AI hardware characteristics.

Scaling to 56 nodes introduced significant cross-node communication overhead, particularly during gradient synchronization in distributed data-parallel training. To mitigate this bottleneck, we leveraged the InfiniBand (IB) networking modules integrated into the Sugon servers. With high-bandwidth, low-latency RDMA support, IB reduces CPU-side data movement and improves the efficiency of cross-node collective communication such as All-Reduce. Together with feature pre-computation for reusable text-side embeddings and image latents, this communication stack helped the 224-device cluster maintain high throughput during both low-resolution pre-training and progressive resolution scaling.

% JuZhou 1.0 represents, to our knowledge, one of the first large-scale text-to-image foundation models whose entire training pipeline---from low-resolution pre-training through progressive resolution scaling to inference step distillation---was completed entirely on domestic Chinese AI accelerators. This experience validates the feasibility of domestic computing infrastructure for frontier generative AI research and provides a practical reference for future software-hardware co-design efforts within the domestic AI ecosystem.
To our knowledge, JuZhou~1.0 is among the first text-to-image foundation models trained and distilled on domestic Chinese AI accelerators, providing practical evidence for large-scale visual generative model training on a domestic software--hardware stack.
